% =============================================================================
% Table 3: Out-of-Sample Forecasting Results
% Fixed train/test split: train Jan 2005 - Dec 2019; test Jan 2025 - Feb 2026.
% 2020-2024 excluded from both windows.
% =============================================================================

\begin{table}[!htbp]
\centering
\small
\caption{Out-of-Sample Forecast Comparison: Mean Squared Prediction Error}
\label{tab:oos_results}

\begin{threeparttable}
\begin{tabular}{l c c c c c}
\toprule
 & \multicolumn{1}{c}{$N$} & \multicolumn{1}{c}{MSPE}
   & \multicolumn{1}{c}{$\%\Delta$ vs.}
   & \multicolumn{1}{c}{DM}
   & \multicolumn{1}{c}{DM} \\
Model
   & \multicolumn{1}{c}{forecasts}
   & \multicolumn{1}{c}{($\times 10^{4}$)}
   & \multicolumn{1}{c}{Baseline}
   & \multicolumn{1}{c}{statistic}
   & \multicolumn{1}{c}{$p$-value} \\
\midrule
Baseline                            & 77 & $0.78$ & ---       & ---     & ---     \\
\addlinespace[2pt]
Homogeneous (HHI + Supp.\ HHI)      & 77 & $0.82$ & $+5.2\%$  & $-2.52$ & $0.012$ \\
Heterogeneous HHI (3A challenger)   & 77 & $0.99$ & $+26.9\%$ & $-1.41$ & $0.159$ \\
Kitchen Sink (3B challenger)        & 77 & $1.12$ & $+43.8\%$ & $-1.58$ & $0.114$ \\
\bottomrule
\end{tabular}

\begin{tablenotes}[flushleft]
\footnotesize
\item \textit{Notes.} Training window: Jan 2005 -- Dec 2019 (1{,}688 observations,
180 months). Test window: Jan 2025 -- Feb 2026 (77 observations, 7 months). The
intervening period (2020--2024) is excluded from both windows. All models include
entity fixed effects only (time fixed effects cannot be projected onto the test
period). The forecast control set is identical across all models: ATM IV, total
return index, term spread (10y minus 2y Treasury), and credit spread (Baa minus
Aaa Moody's seasoned). The challenger models add the structural variables noted
in parentheses. The Diebold--Mariano statistic is computed on monthly-averaged
squared error differentials with $h=1$; a negative DM statistic indicates that
the challenger model has higher MSPE than the baseline. Reading: the structural
challengers do not improve forecast accuracy relative to the baseline; the
Homogeneous specification is in fact significantly worse at the 5\% level. We
interpret this not as a failure of the structural mechanism but as evidence
that the in-sample inference is not the result of overfitting to a pattern
that would also produce gains in monthly out-of-sample point forecasts. See
Section~5.4 for discussion.
\end{tablenotes}
\end{threeparttable}

\end{table}
